\documentclass[10pt]{article}
\usepackage[T1]{fontenc}
\usepackage[utf8]{inputenc}
\usepackage[french]{babel}
\usepackage{amsmath,amssymb,mathrsfs,mathpazo}
\usepackage{geometry}
\geometry{a4paper,margin=2cm}
\usepackage{graphicx,array,tabularx,tabularray,colortbl,multirow,multicol,booktabs,enumitem,pifont,ulem,fancybox,diagbox,cancel}
\usepackage{pstricks,pstricks-add,pst-plot,pst-tree,pst-node,pst-text,pst-func,pst-math,pst-bspline,pst-3dplot,pst-eucl,pst-solides3d}
\let\pscalculate\undefined
\usepackage{pst-bezier}
\usepackage{tikz,pgfplots,tkz-tab}
\usetikzlibrary{arrows,arrows.meta,patterns,calc,positioning,trees,decorations.markings,shapes.geometric,shapes.misc,shapes.arrows,intersections,decorations.pathmorphing,decorations.pathreplacing,backgrounds,shadows,babel}
\pgfplotsset{compat=1.18}
\usepgfplotslibrary{fillbetween}
\usepackage[np]{numprint}
\usepackage{listings,makecell,pas-tableur,twoopt}
\let\barre\undefined
\usepackage[pstricks]{tablvar}
\definecolor{darkgreen}{RGB}{0,100,0}
\providecommand{\red}{\color{red}}
\providecommand{\blue}{\color{blue}}
\providecommand{\green}{\color{green}}
\providecommand{\black}{\color{black}}
\providecommand{\white}{\color{white}}
\providecommand{\gray}{\color{gray}}
\providecommand{\cyan}{\color{cyan}}
\providecommand{\magenta}{\color{magenta}}
\providecommand{\yellow}{\color{yellow}}
\providecommand{\violet}{\color{violet}}
\providecommand{\orange}{\color{orange}}
\providecommand{\euro}{\texteuro}
\providecommand{\EUR}{\texteuro}
\providecommand{\ieme}{\textsuperscript{e}}
\providecommand{\degres}{\ensuremath{^\circ}}
\providecommand{\degree}{\ensuremath{^\circ}}
\providecommand{\vv}[1]{\overrightarrow{#1}}
\providecommand{\Coord}[1]{#1}
\providecommand{\decosix}{}
\providecommand{\starredbullet}{\ensuremath{\bullet}}
\pagestyle{empty}
\setlength{\parindent}{0pt}
\usepackage{adjustbox}\DeclareUnicodeCharacter{2236}{:}\DeclareUnicodeCharacter{2190}{\ensuremath{\leftarrow}}
\providecommand{\R}{\mathbb{R}}
\providecommand{\N}{\mathbb{N}}
\providecommand{\Z}{\mathbb{Z}}
\providecommand{\Q}{\mathbb{Q}}
\providecommand{\D}{\mathbb{D}}
\providecommand{\C}{\mathbb{C}}
\providecommand{\e}{{\mathrm{e}}}
\providecommand{\vect}[1]{\overrightarrow{#1}}
\providecommand{\vectt}[1]{\overrightarrow{\mathrm{#1}}}
\providecommand{\vvt}[1]{\overrightarrow{\mathrm{#1}}}
\providecommand{\np}[2][]{{#2~#1}}
\providecommand{\vv}[1]{{\overrightarrow{#1}}}
\renewcommand{\d}{{\mathrm{d}}}
\providecommand{\degres}{{^\circ}}
\providecommand{\degre}{{^\circ}}
\providecommand{\blue}{\color{blue}}
\providecommand{\red}{\color{red}}
\providecommand{\enskip}{\quad}
\providecommand{\llbracket}{\lbrack\!\lbrack}
\providecommand{\rrbracket}{\rbrack\!\rbrack}
\providecommand{\euro}{{\text{€}}}
\providecommand{\no}{{\text{n°}}}
\providecommand{\interval}[2]{{[#1;#2]}}
\providecommand{\intervFO}[2]{{[#1;#2[}}
\providecommand{\dotfill}{\ldots}

\providecommand{\ds}{\displaystyle}
\providecommand{\textding}[1]{#1}
\providecommand{\Oij}{\(\mathrm{O};\vec{\imath},\vec{\jmath}\)}
\providecommand{\Oijk}{\(\mathrm{O};\vec{\imath},\vec{\jmath},\vec{k}\)}
\providecommand{\Ouv}{\(\mathrm{O};\vec{u},\vec{v}\)}
\providecommand{\intervFF}[2]{\ensuremath{[#1;#2]}}
\renewcommand{\euro}{\ifmmode\text{\texteuro}\else\texteuro\fi}
\DeclareUnicodeCharacter{00B0}{\ifmmode^{\circ}\else\textdegree\fi}
\DeclareUnicodeCharacter{20AC}{\ifmmode\text{\texteuro}\else\texteuro\fi}
% Énoncés : Éducation nationale. Transcription et corrigés : APMEP, Denis Vergès et contributeurs. Découpage et classement : CoopMaths, Sébastien Lozano, Éric Elter et contributeurs. Adaptation de la présentation : Maths Explorer.
% CC BY-SA : https://coopmaths.fr/www/about/
% https://coopmaths.fr/alea/static/bac/2024/tex/bac_2024_06_sujet1_asie_2.tex
\begin{document}
\noindent\textit{APMEP / CoopMaths. CC BY-SA. Présentation adaptée par Maths Explorer.}\par\medskip
\medskip

\begin{minipage}{0.48\linewidth}
Dans l'espace muni d'un repère orthonormé \Oijk{} d'unité 1 cm, on considère les points : A$(3~;~-1~;~1)$ ; B$(4~;~-1~;~0)$ ;
 C$(0 ~;~3~;~2)$ ; D$(4~;~3~;~-2)$ et S$(2~;~1~;~4)$.

Dans cet exercice on souhaite montrer que SABDC est une pyramide à base ABDC trapézoïdale de sommet S, afin de calculer son volume.
\end{minipage}\hfill
\begin{minipage}{0.48\linewidth}
\psset{unit=1cm}
\begin{pspicture}(4.7,5)
\pspolygon(2.4,4.5)(0.2,1.5)(1.3,0.2)(4.3,2.7)(2.4,4.5)(1.3,0.2)
\psline[linestyle=dashed](2.4,4.5)(1.4,2.6)(0.2,1.5)
\psline[linestyle=dashed](1.4,2.6)(4.3,2.7)
\uput[d](1.4,2.6){A} \uput[ul](0.2,1.5){B} \uput[dr](4.3,2.7){C} \uput[d](1.3,0.2){D} \uput[u](2.4,4.5){S}
\end{pspicture}
\end{minipage}

\medskip

\begin{enumerate}
\item Montrer que les points A, B et C ne sont pas alignés.
\item
	\begin{enumerate}
		\item Montrer que les points A, B, C et D sont coplanaires.
		\item Montrer que le quadrilatère ABDC est un trapèze de bases [AB] et [CD].

\emph{On rappelle qu'un trapèze est un quadrilatère ayant deux côtés opposés parallèles appelés bases}.
	\end{enumerate}
\item
	\begin{enumerate}
		\item Démontrer que le vecteur $\vect{n}(2~;~1~;~2)$ est un vecteur normal au plan (ABC).
		\item En déduire une équation cartésienne du plan (ABC).
		\item Déterminer une représentation paramétrique de la droite $\Delta$ passant par le point S et orthogonale au plan (ABC).
		\item On note I le point d'intersection de la droite $\Delta$ et du plan (ABC).

Montrer que le point I a pour coordonnées $\left(\dfrac23~;~\dfrac13~;~\dfrac83~\right)$, puis montrer que $\mathrm{SI} = \np[cm]{2}$.
		\end{enumerate}
\item
	\begin{enumerate}
		\item Vérifier que le projeté orthogonal H du point B sur la droite (CD) a pour coordonnées H$(3~;~3~;~-1)$ et montrer que HB $= 3\sqrt 2$~cm.
		\item Calculer la valeur exacte de l'aire du trapèze ABDC.

On rappelle que l'aire d'un trapèze est donnée par la formule
\[\mathcal{A} = \dfrac{b + B}{2} \times h\]

où $b$ et $B$ sont les longueurs des bases du trapèze et $h$ sa hauteur.
		\end{enumerate}
\item Déterminer le volume de la pyramide SABDC.

On rappelle que le volume $V$ d'une pyramide est donné par la formule

\begin{center}$V = \dfrac13 \times$  aire de la base $\times$ hauteur\end{center}
\end{enumerate}

\bigskip

\end{document}
